\section{A paired-coordinate cube application}\label{h:cube:section}

The exact-density theorem for finite measures also applies to a prior
supported on a lower-dimensional subset of the sphere.  We illustrate
this with a cube embedded by coordinate pairs.  Its first coordinates
transfer the local mass bound to ambient Euclidean balls, and its
Lipschitz bound supplies a common support for the projected labels.
These are the hypotheses needed by the general projection and selection
estimates.

The resulting success bound concerns the cube prior.  It gives a
sample lower bound at fine accuracy from cube-average success alone.
To cover the remaining accuracies for a learner that succeeds at every
unit signal, we use the uniform-sphere endpoint from the main proof.

\subsection{The prior and its terminal mass}

Let \(d\ge16\) be an integer, set \(n=\lfloor d/2\rfloor\), and give
\[
 Q=[-1/2,1/2]^n
\]
Lebesgue probability measure \(du\).  The cube has volume one.  Define
\begin{equation}\label{h:cube:map}
 \iota(u)=\frac1{\sqrt n}
 \bigl(u_1,\ldots,u_n,\sqrt{1-u_1^2},\ldots,
       \sqrt{1-u_n^2},0,\ldots,0\bigr)\in\mathbb R^d,
\end{equation}
with \(d-2n\) zero coordinates.  Each coordinate pair contributes
\(1/n\) to the squared norm, so \(\iota(Q)\subset S^{d-1}\).
The signal prior is \(\iota_\#du\).

Write \(\operatorname{pr}:\mathbb R^d\to\mathbb R^n\) for projection
onto the first \(n\) coordinates.  Then
\(\operatorname{pr}\iota(u)=u/\sqrt n\).  On \([-1/2,1/2]\), the
derivative of \(a\mapsto\sqrt{1-a^2}\) has absolute value at most
\(1/\sqrt3\).  The mean value theorem therefore gives
\begin{equation}\label{h:cube:lipschitz}
 \|\iota(u)-\iota(v)\|\le\frac2{\sqrt n}\|u-v\|
 \qquad(u,v\in Q).
\end{equation}

Set
\begin{equation}\label{h:cube:parameters}
 \alpha=\frac n4,\qquad \beta=n-\alpha=\frac{3n}{4},
 \qquad p=\left\lfloor\frac n4\right\rfloor\ge2.
\end{equation}
For \(z\in\mathbb R^n\) and \(r>0\), abbreviate
\[
 \mathcal B(z,r)=Q\cap B_n(z,\sqrt n\,r),
\]
where \(B_j(x,t)\) is the closed Euclidean ball in \(\mathbb R^j\).
The local bound for a function \(f:Q\to[0,1]\) will be
\begin{equation}\label{h:cube:local}
 \int_{\mathcal B(z,r)}f(u)\,du\le Lr^\beta
 \qquad(z\in\mathbb R^n,\ r>0).
\end{equation}
We will propagate the coefficient \(L\) backward through the stream.

Fix a terminal state and stopping index, after fixing the independent
rule seed.  Its output law is independent of the signal.  Let
\(f_{\rm term}(u)\) be the probability that this output has angular
error at most \(\epsilon\) at \(\iota(u)\).  For a fixed unit output
\(\widehat s\), angular error at most \(\epsilon\) implies
chordal error at most \(\epsilon\).  Projection onto the first
coordinates then gives
\[
 \|u-\sqrt n\,\operatorname{pr}\widehat s\|
 \le\sqrt n\,\epsilon.
\]
Let \(\omega_j\) be the volume of the unit ball in \(\mathbb R^j\).
The Gaussian integral gives
\(\omega_j j^{j/2}\le(2\pi e)^{j/2}\): integrate
\(e^{-\|x\|^2/2}\) over the ball of radius \(\sqrt j\) and compare
with its integral over \(\mathbb R^j\).  The successful part of
\(\mathcal B(z,r)\) lies in two parameter balls, so its volume is at
most the smaller of their volumes.  With \(C_0=\sqrt{2\pi e}\), the
terminal success function consequently satisfies
\begin{equation}\label{h:cube:terminal}
 \int_{\mathcal B(z,r)}f_{\rm term}(u)\,du
 \le C_0^n\min(r,\epsilon)^n
 \le C_0^n\epsilon^\alpha r^\beta.
\end{equation}
The last inequality uses \(\alpha+\beta=n\); the same estimate for a
random output follows by integration against its signal-independent law.

\subsection{Transferring local mass to the embedded measure}

Let \(\gamma_k\) be the law of a \(k\)-by-\(d\) matrix with
independent standard Gaussian entries, and put
\[
 \lambda_k(dA,dy)=\gamma_k(dA)\,dy.
\]
The following specialization supplies both the exact density and its
common support.  A parameter function measurable in the Lebesgue
completion can first be replaced by a Borel version under \(du\),
without changing any local integral.

\begin{lemma}[Exact density and a common projected ellipsoid]
\label{h:cube:density}
Let \(k\) be an integer with \(1\le k\le p\), and let a Borel function
\(f:Q\to[0,1]\) satisfy \eqref{h:cube:local} with \(L\ge0\).
For \(z\in\mathbb R^n\) and \(r>0\), define the finite joint measure
\begin{equation}\label{h:cube:joint-measure}
 \eta_{f,z,r}(E)
 =\int\!\int_{\mathcal B(z,r)}
       \mathbf 1_E(A,A\iota(u))f(u)\,du\,\gamma_k(dA)
\end{equation}
for Borel \(E\subseteq\mathbb R^{k\times d}\times\mathbb R^k\).
There is a universal \(C_2\ge1\) such that
\(\eta_{f,z,r}=h_{f,z,r}\lambda_k\) for a nonnegative density with
\begin{equation}\label{h:cube:density-bound}
 \|h_{f,z,r}\|_{L^p(\lambda_k)}
 \le C_2^n Lr^{\beta-k(1-1/p)}.
\end{equation}
If \(\mathcal B(z,r)\ne\varnothing\), fix any
\(u_*\in\mathcal B(z,r)\) and set
\begin{equation}\label{h:cube:ellipsoid-support}
 \mathcal E_{u_*,r}
 =\{(A,y):y\in A\iota(u_*)+A B_d(0,4r)\}.
\end{equation}
Every such density, for this same \(z,r,u_*\), vanishes almost
everywhere outside \(\mathcal E_{u_*,r}\), and
\begin{equation}\label{h:cube:ellipsoid-volume}
 \lambda_k(\mathcal E_{u_*,r})\le C_2^n r^k.
\end{equation}
\end{lemma}

\begin{proof}
An empty \(\mathcal B(z,r)\) gives the zero measure.  Otherwise fix
\(u_*\in\mathcal B(z,r)\), and let
\[
 \mu=\iota_\#\bigl(\mathbf1_{\mathcal B(z,r)}f(u)\,du\bigr).
\]
For any \(x\in\mathbb R^d\) and \(t>0\), projection onto the first
coordinates gives
\begin{equation}\label{h:cube:ambient-growth}
 \begin{split}
 \mu(B_d(x,t))
 &\le\int_{\mathcal B(\sqrt n\,\operatorname{pr}x,t)}f(u)\,du\\
 &\le Lt^\beta.
 \end{split}
\end{equation}
For \(u\in\mathcal B(z,r)\), the distance from \(u\) to \(u_*\) is
at most \(2\sqrt n\,r\).  Thus \eqref{h:cube:lipschitz} puts
\(\mu\) inside \(B_d(\iota(u_*),4r)\).

If \(L=0\), \eqref{h:cube:local} applied to a ball containing \(Q\)
gives \(f=0\) almost everywhere, and the zero density suffices.  Suppose
now that \(L>0\).  Apply the general finite-measure result
Theorem~\ref{h:dv:finite-density} with
\[
 H=L,\qquad R=4r,\qquad m=k,
 \qquad \beta=3n/4,\qquad p=\lfloor n/4\rfloor.
\]
For its factors \(J_j\), the required margin and their uniform bounds
are
\begin{equation}\label{h:cube:margin}
 \beta-j-k\ge\frac{3n}{4}-2p+2\ge\frac n4+2,
 \qquad
 J_j=1+3^j2^\beta\frac{k}{\beta-j-k}
 \le1+3^{n/4}2^{3n/4}
 \qquad(0\le j\le p-2).
\end{equation}
In particular \(\beta>k+p-2\).  We use the general parameter range
of that theorem, not the specialization to exponents near half
the ambient dimension.
Its exact norm bound is
\[
 (2\pi)^{-k(p-1)/(2p)}
 L(4r)^{\beta-k(1-1/p)}
 \left(\prod_{j=0}^{p-2}J_j\right)^{1/p}.
\]
The first factor is at most one, each \(J_j\le C^n\), and
\(0<\beta-k(1-1/p)\le\beta\).  Absorbing
\(4^{\beta-k(1-1/p)}\le4^\beta\) into \(C_2^n\) proves
\eqref{h:cube:density-bound}.  This is the density of the exact joint
measure \eqref{h:cube:joint-measure}; no density relative to spherical
area is assumed.

The set \(\mathcal E_{u_*,r}\) is the exact projected image of the
enclosing ball \(B_d(\iota(u_*),4r)\).  Lemma~\ref{h:ev:ellipsoid}
makes this set Borel and gives the asserted common support.  The same
lemma gives
\[
 \lambda_k(\mathcal E_{u_*,r})
 \le C_{\mathrm E}^d(4r)^k\le C_2^n r^k,
\]
after one universal increase of \(C_2\), since
\(d\le2n+1\) and \(k\le n/4\).  This calculation also covers the
zero-density case.
\end{proof}

\subsection{The block bound and the streaming consequence}

The power of \(r\) lost in the density norm is restored by the common
support volume.  Selection then costs only the \(p\)th root of the
number of continuations.

\begin{corollary}[A block on the paired-coordinate prior]
\label{h:cube:block}
Let \(k,N\) be integers with \(1\le k\le p\) and \(N\ge1\), and let Borel
\(f_1,\ldots,f_N:Q\to[0,1]\) satisfy \eqref{h:cube:local} with the
same \(L\ge0\).  Suppose Borel functions
\(g_v:\mathbb R^{k\times d}\times\mathbb R^k\to[0,\infty)\) obey
\(\sum_{v=1}^N g_v(A,y)\le1\).  Then
\[
 F(u)=\mathbb E_{A\sim\gamma_k}
       \sum_{v=1}^N g_v(A,A\iota(u))f_v(u)
\]
satisfies \eqref{h:cube:local} with \(C_3^nN^{1/p}L\) in place of
\(L\), for a universal \(C_3\ge1\).
\end{corollary}

\begin{proof}
Fix \(z,r\).  An empty \(\mathcal B(z,r)\) is immediate.  Otherwise
choose the same \(u_*\in\mathcal B(z,r)\) for all \(v\).  The exact
joint densities from Lemma~\ref{h:cube:density} share its region
\(\mathcal E_{u_*,r}\).  Apply
Lemma~\ref{h:alt:selection} with density norm at most
\(C_2^nLr^{\beta-k(1-1/p)}\) and domain measure at most
\(C_2^nr^k\).  The image-measure identity and that lemma give
\begin{equation}\label{h:cube:selection}
 \int_{\mathcal B(z,r)}F(u)\,du
 \le N^{1/p}C_2^n(C_2^n)^{1-1/p}Lr^\beta
 \le C_3^nN^{1/p}Lr^\beta.
\end{equation}
The radius powers cancel, and the same \(C_3\) works for every
\(1\le k\le p\).
\end{proof}

To apply this bound to the model's completed-measurable rules, first
observe that \(f=1\) satisfies
\[
 |\mathcal B(z,r)|\le\min\{1,C_0^nr^n\}\le C_0^nr^\beta.
\]
Take \(k=1\) in Lemma~\ref{h:cube:density} and a parameter ball
containing \(Q\).  The fresh pair
\((X,\langle X,\iota(u)\rangle)\), for \(u\sim du\), is therefore
absolutely continuous relative to \(\gamma_1(dx)\,dy\).
We may consequently use the seedwise Borel reduction and fresh-state
continuations of Section~\ref{h:stream:reduction} under this cube prior.
As there, the conditional numerical estimates are averaged using the
original measurable experiment; no jointly chosen family of Borel
replacements is required.

\Needspace{4\baselineskip}
\begin{proposition}[Precision bound for the cube prior]
\label{h:cube:success}
There are universal constants \(D\ge e\), \(d_0\ge16\), and \(c>0\)
with the following property.  Let \(d\ge d_0\), \(M\ge0\), and \(T\ge0\)
be integers, with \(M\le d^2\) and \(0<\epsilon\le1/10\).
Every learner in Definition~\ref{h:model} with these memory and sample
bounds has angular success probability under \(s=\iota(u)\), \(u\sim du\),
at most
\begin{equation}\label{h:cube:success-bound}
 \min\left\{1,\,
       D^{n(m+1)}\epsilon^{n/4}\right\},
 \qquad
 m=\left\lceil\frac{T}{k_0}\right\rceil,\qquad
 k_0=\left\lfloor\frac n4\right\rfloor.
\end{equation}
The probability averages over \(u\), the fresh Gaussian samples, and
the learner's randomness.  If that same learner has angular success
probability at least \(2/3\) for every \(s\in S^{d-1}\), then
\begin{equation}\label{h:cube:sample-result}
 T\ge c\,d\log(1/\epsilon).
\end{equation}
\end{proposition}

\begin{proof}
Use the seedwise convention just established, and first start from any
prescribed initial state.  For \(T>0\), partition the stream into
\(m=\lceil T/p\rceil\) consecutive blocks, each of length at most
\(p=k_0\).  Lemma~\ref{h:stream:categories}, with angular terminal
score and signal \(\iota(u)\), represents the preceding success
function by the exact block formula in Corollary~\ref{h:cube:block}.
For a block of length \(k\), its Borel routes have at most
\begin{equation}\label{h:cube:destinations}
 N\le(k+2)2^M
\end{equation}
destinations.  The stopping offsets are local to this block, including
an immediate stop; its continuing functions use fresh future samples.
These are the same phase and terminal-state conventions as in that
lemma.

For \(n\ge8\), \(p\ge n/8\), \(d\le2n+1\), and \(M\le d^2\),
\[
 \frac{\log N}{p}
 \le\frac{d^2\log2+\log(p+2)}p\le C_4n
\]
with a universal \(C_4\).  Choose a universal \(D\ge e\) so that
\(D\ge C_0\) and \(C_3^nN^{1/p}\le D^n\).
The terminal coefficient is \(C_0^n\epsilon^{n/4}\) by
\eqref{h:cube:terminal}.  Backward induction, applying the block bound
and retaining the terminal estimate at earlier stops, gives coefficient
\[
 L=D^{n(m+1)}\epsilon^{n/4}
\]
for every initial state.  The ball \(\mathcal B(0,1)\) contains \(Q\),
so this coefficient bounds total cube-average success.  For \(T=0\),
the terminal bound gives the same statement with \(m=0\).
All constants are uniform over the seedwise Borel rules and prescribed
states.  The independent conditional initialization and the averaging
of original conditional success probabilities from
Section~\ref{h:stream:reduction} therefore prove
\eqref{h:cube:success-bound} for the original learner, including every
short last block.

Suppose now only that the cube-average success is at least \(2/3\).
Taking natural logarithms in \eqref{h:cube:success-bound} and using
\(m\le T/k_0+1\) gives
\begin{equation}\label{h:cube:sample-bound}
 \frac{T}{k_0}
 \ge\frac{\log(1/\epsilon)}{4\log D}
       -2-\frac{\log(3/2)}{n\log D}.
\end{equation}
If \(\log(1/\epsilon)\ge24\log D\), its right side is at least
\(\log(1/\epsilon)/(8\log D)\).  Since \(k_0\ge n/8\) and
\(n\ge d/3\), this already gives
\(T\ge d\log(1/\epsilon)/(192\log D)\) in that accuracy range.

For the all-accuracy conclusion \eqref{h:cube:sample-result}, assume
the stated every-signal guarantee.  It gives the cube-average premise
just used.  It also gives uniform-sphere average success at least
\(2/3\) for the same original learner.  The reflection consequence
\eqref{h:end:linear-threshold} therefore forces \(T>d/4\).
For \(\log(1/\epsilon)<24\log D\), this implies
\[
 T>\frac d4>\frac{d\log(1/\epsilon)}{96\log D}.
\]
Combining the two ranges proves \eqref{h:cube:sample-result}, for
example with \(c=1/(192\log D)\), after increasing the universal
dimension threshold as needed.  The memory condition \(M\le d^2\)
in particular holds eventually for every sequence \(M(d)=o(d^2)\).
\end{proof}

\subsection{An affine-distance interpretation}

The application is complete.  The embedding also gives a direct
interpretation of the Gram moment behind its projection bound.  This
alternative calculation isolates what the first coordinates guarantee
about affine spans.  For \(u,u_1,\ldots,u_j\in Q\), \(j\ge1\),
\begin{equation}\label{h:cube:affine}
 \operatorname{dist}\bigl(\iota(u),
       \operatorname{aff}(\iota(u_1),\ldots,\iota(u_j))\bigr)
 \ge\frac1{\sqrt n}
       \operatorname{dist}\bigl(u,\operatorname{aff}(u_1,\ldots,u_j)\bigr).
\end{equation}
Indeed, for \(\sum_i c_i=1\), projection of
\(\iota(u)-\sum_i c_i\iota(u_i)\) is
\((u-\sum_i c_i u_i)/\sqrt n\).  Projection cannot increase norm,
and taking the infimum proves the claim, including dependent points.

\begin{lemma}[Affine inverse moments]\label{h:cube:inverse}
Let \(f:Q\to[0,1]\) be Borel and satisfy \eqref{h:cube:local} with
\(L\ge0\).  Let \(k\) be an integer with \(1\le k\le p\), and let
\(V\subseteq\mathbb R^n\) be a nonempty affine subspace of dimension
\(j\le p-2\).  There is a universal \(C_1\ge1\) such that, for every
\(z\in\mathbb R^n\), \(r>0\), and \(0<t\le1\),
\begin{align}
 \int_{\mathcal B(z,r)}f(u)
   \mathbf1_{\{\operatorname{dist}(u,V)\le\sqrt n\,tr\}}\,du
 &\le C_1^n Lr^\beta t^{\beta-j},
 \label{h:cube:tube}\\
 \int_{\mathcal B(z,r)}f(u)
   \left(\frac{\operatorname{dist}(u,V)}{\sqrt n}\right)^{-k}\,du
 &\le C_1^n Lr^{\beta-k}.
 \label{h:cube:inverse-bound}
\end{align}
The second integrand is nonnegative and extended-valued on \(V\).
\end{lemma}

\begin{proof}
The cases \(L=0\) and an empty parameter ball are immediate as above.
Apply Proposition~\ref{h:dv:ambient-gram} in parameter dimension
\(n\) to
\(\mathbf1_{\mathcal B(z,r)}f(u)\,du\), with
\[
 H=Ln^{-\beta/2},\qquad R=\sqrt n\,r,\qquad m=k.
\]
The growth hypothesis follows because a parameter ball of radius
\(a>0\) is \(\mathcal B(x,a/\sqrt n)\) after intersection with
\(Q\).  Its tube bound, evaluated at \(\sqrt n\,tr\), is
\[
 3^j2^\beta (Ln^{-\beta/2})(\sqrt n\,r)^j
                 (\sqrt n\,tr)^{\beta-j}
 =3^j2^\beta Lr^\beta t^{\beta-j}.
\]
Its inverse-distance bound is
\(J_jLn^{-k/2}r^{\beta-k}\).  Multiplication by \(n^{k/2}\)
gives the normalization in \eqref{h:cube:inverse-bound}.
The bounds \eqref{h:cube:margin} absorb both constants into a universal
\(C_1^n\).
\end{proof}

For example, let \(G\) be the Gram matrix of the \(p-1\) differences
\(\iota(u_i)-\iota(u_1)\), \(2\le i\le p\).
Almost every parameter tuple is affinely independent, since proper
affine subspaces of \(\mathbb R^n\) have zero Lebesgue measure.
Writing \(V_{i-1}=\operatorname{aff}(u_1,\ldots,u_{i-1})\),
Gram--Schmidt and \eqref{h:cube:affine} give
\begin{equation}\label{h:cube:gram}
 \det G\ge
 \prod_{i=2}^p
 \left(\frac{\operatorname{dist}(u_i,V_{i-1})}{\sqrt n}\right)^2>0.
\end{equation}
Integrating in the order \(u_p,\ldots,u_2\), followed by \(u_1\),
therefore yields
\[
 \int_{\mathcal B(z,r)^p}\det G^{-k/2}
             \prod_{i=1}^p f(u_i)\,du_1\cdots du_p
 \le C_1^{n(p-1)}L^p r^{\beta p-k(p-1)}.
\]
Thus one can also obtain the projection bound by inserting this Gram
estimate into the Gaussian difference identity and compact-test passage
of Section~\ref{h:projection}.  The main application instead uses the
general finite-measure theorem through the ambient ball bound
\eqref{h:cube:ambient-growth}.

\paragraph{An enclosing-ball alternative.}
The common ellipsoid may instead be enclosed in the larger label ball
\begin{equation}\label{h:cube:support}
 \mathcal R_A=B_k\bigl(A\iota(u_*),4r\|A\|_{\rm op}\bigr).
\end{equation}
Thus every density in Lemma~\ref{h:cube:density}, with the same
\(z,r,u_*\), also vanishes almost everywhere outside
\(\{(A,y):y\in\mathcal R_A\}\).  Its volume has the bound
\begin{equation}\label{h:cube:volume}
 \mathbb E_{A\sim\gamma_k}|\mathcal R_A|
 =(4r)^k\mathbb E\,\vol_k B_k(0,\|A\|_{\rm op})
 \le4^k e^{Cd}r^k\le C_2^n r^k,
\end{equation}
after a universal increase of \(C_2\), where
\(|\mathcal R_A|\) is \(k\)-dimensional Lebesgue volume.
This estimate follows independently from the Gaussian net bound in
Lemma~\ref{h:proj:volume}, not from the smaller ellipsoid's volume.
It gives the same block bound, although the direct proof above uses
the exact projected ellipsoid and does not require this enlargement.
